\subsection{Población y Muestra}
Dado que este estudio utiliza un enfoque cuantitativo basado en
mediciones numéricas, la recolección de datos se apoya en valores
reales que representen fielmente la situación de la clínica. Para
garantizar representatividad y evitar sesgos, los ciclos de atención
administrativa que conforman la segunda unidad de análisis serán
seleccionados mediante muestreo aleatorio simple. Este procedimiento
consiste en asignar un número de control consecutivo a cada ciclo
registrado durante el período de estudio y elegir aleatoriamente
los que formarán parte de la medición en cada fase. De este modo,
cada ciclo de atención tiene la misma probabilidad de ser incluido,
lo que permite comparar los resultados antes y después de la
implementación del sistema con validez estadística.

La primera unidad de análisis está conformada por el personal
directamente involucrado en la atención y administración del centro
odontológico Bellidental. Al tratarse de un establecimiento de
pequeña escala, el número total de personas que participan en la
gestión es reducido y perfectamente conocido; por ello se trabajará
con un censo total, incluyendo a todos sus miembros sin necesidad
de seleccionar una muestra parcial. Este grupo constituye la fuente
principal de información para comprender los procesos actuales, será
el insumo central para el diseño del nuevo sistema de información y
el sujeto de observación durante las mediciones del ciclo de
atención. El personal del proceso de atención y administración se
detalla en la Tabla~\ref{tab:actores_proceso}.

\begin{table}[H]
\centering
\begin{tabular}{|l|c|p{7.5cm}|}
\hline
\textbf{Rol} & \textbf{Cantidad} &
\textbf{Participación en el estudio} \\ \hline
Odontólogo general      & 1 & Entrevista, revisión del flujo de
                              trabajo y usuario final del sistema \\ \hline
Odontólogo especialista & 1 & Entrevista, revisión del flujo de
                              trabajo y usuario final del sistema \\ \hline
Secretaria              & 1 & Entrevista, revisión de tareas de
                              oficina y usuario final del sistema \\ \hline
\textbf{Total}          & \textbf{3} & \\ \hline
\end{tabular}
\caption{Personal del proceso de atención y administración del
         centro odontológico Bellidental.}
\label{tab:actores_proceso}
\end{table}

La segunda unidad de análisis corresponde a los ciclos completos de
atención administrativa. Siguiendo la definición operacional
establecida en la delimitación del estudio, cada ciclo se define
como el período transcurrido desde que el paciente es registrado en
recepción hasta que se completa el asiento del pago y se emite el
comprobante correspondiente, abarcando los subprocesos de registro
de datos personales, apertura o actualización de la historia
clínica, consulta odontológica, elaboración de receta y registro
de facturación. Dado que la clínica mantiene un flujo regular de
aproximadamente 75 pacientes semanales, el volumen total de
atenciones durante el período de estudio es limitado y cuantificable.
El número exacto de ciclos a observar en cada fase se determinará
aplicando la fórmula para poblaciones finitas con variable continua,
considerando un nivel de confianza del 95\,\%, una estimación de la
desviación estándar obtenida a partir de los tiempos registrados en
estudios análogos de entornos dentales similares~\cite{Ozkaynak2021},
y un margen de error preestablecido de dos minutos, congruente con la
precisión requerida para detectar diferencias clínicamente
relevantes en el ciclo de atención. El total resultante se dividirá
en dos grupos de igual tamaño, uno correspondiente a la fase previa
a la implementación y otro a la fase posterior, para calcular el
tiempo promedio y la dispersión en cada etapa y evaluar si la
diferencia observada es estadísticamente significativa. Las fases
de medición se presentan en la Tabla~\ref{tab:fases_medicion}.

\begin{table}[H]
\centering
\begin{tabular}{|p{3.2cm}|p{3.8cm}|p{6.2cm}|}
\hline
\textbf{Fase de medición} & \textbf{Período} &
\textbf{Criterio de selección} \\ \hline
Pre-implementación  & Julio--agosto de 2026
                    & Muestreo aleatorio simple sobre
                      los ciclos completos registrados
                      durante la fase de diagnóstico \\ \hline
Post-implementación & Noviembre de 2026
                    & Muestreo aleatorio simple sobre
                      los ciclos completos registrados
                      con el sistema en operación \\ \hline
\end{tabular}
\caption{Fases de medición de ciclos de atención administrativa.}
\label{tab:fases_medicion}
\end{table}

Los criterios de inclusión considerarán únicamente los ciclos en
los que el paciente complete las tres etapas del proceso en la
misma visita: registro, consulta y cobro. Se excluirán los casos
de urgencias no programadas, las consultas que no generen receta
ni registro de pago, y los ciclos interrumpidos por causas externas
ajenas al flujo administrativo del consultorio.
